\chapter{Grafy plumbingowe, formy przecięcia i krata \texorpdfstring{$E_8$}{E8}}
\label{ch:tg-plumbing}

W poprzednich rozdziałach rozpatrywaliśmy dwie konstrukcje: z jednej
strony grafy, których wierzchołki kodują części rozmaitości, z drugiej
strony ramowane sploty opisujące dołączanie uchwytów indeksu dwa.
\emph{Plumbing} łączy te punkty widzenia. Z prostego grafu z wagami
całkowitymi budujemy gładką czterorozmaitość $P_\Gamma$, a następnie
odczytujemy z grafu jej formę przecięcia i homologię brzegu.
To jeden z najbardziej użytecznych sposobów przechodzenia od
kombinatoryki do topologii czterowymiarowej.

Zaczniemy od drzew o wierzchołkach reprezentujących sfery. W tym
przypadku konstrukcja jest elementarna i pozwala przeprowadzić
rachunki bez rozbudowanego aparatu rozmaitości Seiferta. Na końcu
omówimy ograniczenia metody oraz jej związek z teorią Donaldsona
i klasyfikacją Freedmana. Wszystkie rozmaitości plumbingowe
rozważamy jako \emph{zwarte rozmaitości gładkie z brzegiem};
nie są one automatycznie rozmaitościami zamkniętymi.

\section{Wiązki dyskowe i lokalna operacja plumbing}
\label{sec:tg-plumbing-definition}

\begin{definicja}[Zorientowana wiązka dysków nad sferą]
Niech $e\in\mathbb Z$. Przez $E(e)\to S^2$ oznaczamy zorientowaną
wiązkę $D^2$ o grupie strukturalnej $SO(2)$, której skojarzona
zorientowana wiązka wektorowa rzędu dwa ma klasę Eulera
$e\,u\in H^2(S^2;\mathbb Z)$, gdzie $u$ jest generatorem
wyznaczonym przez orientację sfery. \emph{Sekcją zerową}
nazywamy obraz odwzorowania $s\colon S^2\to E(e)$ przypisującego
punktowi bazowemu wektor zerowy jego włókna.
\end{definicja}

Każdą zorientowaną wiązkę płaszczyzn nad $S^2$ klasyfikuje
liczba całkowita: stopień funkcji przejścia
$S^1\to SO(2)\cong S^1$ między dwiema trywializacjami
na półkulach. Jest to właśnie liczba Eulera $e$.
Jeżeli przez $S$ oznaczymy obraz sekcji zerowej,
to jej wiązka normalna w $E(e)$ jest pierwotną wiązką
płaszczyzn, zatem $S\cdot S=e$. Równość tę należy
rozumieć jako liczbę przecięcia $S$ z niewielkim,
generycznym przesunięciem w kierunku normalnym.
Dokładniej, wybieramy sekcję normalną poprzeczną
do sekcji zerowej. Jej zera odpowiadają przecięciom
obu sekcji, a lokalne indeksy zer są znakami
przecięć przy orientacji iloczynowej bazy i włókna.
Suma tych indeksów jest równa
$\langle e(E),[S^2]\rangle=e$.
Odwrócenie orientacji samej reprezentującej klasy
$[S]$ nie zmienia jej samoprzecięcia.

\begin{definicja}[Dodatni plumbing dwóch wiązek]
Niech $E_1\to S^2$ i $E_2\to S^2$ będą zorientowanymi
wiązkami dyskowymi. W każdej wybieramy dysk $D_i\subset S^2$
nieprzecinający dysków wybranych dla innych operacji.
Nad $D_i$ ustalamy trywializację, która daje blok
$D^2_{\text{baza}}\times D^2_{\text{włókno}}$.
W lokalnym modelu plumbingowania identyfikujemy wybrane
bloki obu wiązek (po odpowiednim ograniczeniu ich
trywializacji), używając odwzorowania
\[
(z,w)\longmapsto(w,z), \qquad (z,w)\in D^2\times D^2,
\]
z dobranymi orientacjami, tak aby sekcje zerowe przecinały się
w powstałym bloku dodatnio i poprzecznie w jednym punkcie.
Wybieramy dyski bazowe we wnętrzach większych
dysków trywializujących. W ilorazie dwóch wiązek
ściana bazowa pierwszego bloku przechodzi na
ścianę włóknową drugiego i odwrotnie. Wzdłuż
wnętrza ściany bazowej jedna wiązka dostarcza
przedłużenia kołnierza, a druga kończy się na
ścianie włóknowej; nie powstaje rozgałęzienie
otoczenia. Przy przecięciu ścian otrzymujemy
lokalny model narożnika wklęsłego, który
wygładzamy. Jest to standardowa konstrukcja
plumbingu wiązek dyskowych (por. R.~E. Gompf,
A.~I. Stipsicz, \emph{4-Manifolds and Kirby Calculus},
AMS, 1999, sekcja 4.6).
Po wygładzeniu narożników otrzymujemy \emph{dodatni plumbing}.
\end{definicja}

\begin{uwaga}[Lokalna geometria i znaki]
W bloku $D^2\times D^2$ sekcje zerowe dwóch wiązek
reprezentują odpowiednio dyski $D^2\times\{0\}$ i
$\{0\}\times D^2$. Mają jeden punkt przecięcia $(0,0)$.
Zamiana czynników zachowuje orientację czterowymiarową,
ponieważ permutuje dwa czynniki rzeczywiste wymiaru dwa.
Orientacje samych sfer i wiązek trzeba dobrać konsekwentnie:
wówczas znak lokalnego przecięcia wynosi $+1$. Zmiana
orientacji jednej sfery odwraca znaki wszystkich jej przecięć
z sąsiednimi sferami, ale nie zmienia jej samoprzecięcia.
\end{uwaga}

\begin{center}
\begin{tikzpicture}[scale=.98,line cap=round,>=Stealth]
  \begin{scope}[shift={(-2.65,0)}]
    \fill[tgBlue!15] (-1,-.65) rectangle (1,.65);
    \draw[tgBlue,very thick] (-1,-.65) rectangle (1,.65);
    \fill[tgCoral!22] (-.27,-.65) rectangle (.27,.65);
    \draw[tgCoral,line width=2.6pt] (0,-.65)--(0,.65);
    \draw[tgBlue,line width=2.4pt] (-1,0)--(1,0);
    \node[text=tgNavy,font=\footnotesize] at (0,-1.05) {$D^2_{\text{baza}}\times D^2_{\text{włókno}}$};
  \end{scope}
  \draw[tgTeal,-{Stealth[length=3mm]},line width=1.6pt] (-.75,.05)--(.75,.05);
  \node[text=tgTeal,font=\footnotesize] at (0,.42) {zamiana ról};
  \begin{scope}[shift={(2.65,0)}]
    \fill[tgCoral!15] (-1,-.65) rectangle (1,.65);
    \draw[tgCoral,very thick] (-1,-.65) rectangle (1,.65);
    \draw[tgBlue,line width=2.4pt] (0,-.65)--(0,.65);
    \draw[tgCoral,line width=2.6pt] (-1,0)--(1,0);
    \fill[tgPurple] (0,0) circle (.09);
    \node[text=tgNavy,font=\footnotesize] at (0,-1.05) {przecięcie dwóch sekcji};
  \end{scope}
  \node[text=tgNavy,font=\scriptsize] at (0,-1.52)
    {Model lokalny: jeden punkt przecięcia, bez utożsamiania całych sfer.};
\end{tikzpicture}
\end{center}

\begin{definicja}[Ważone drzewo plumbingowe]
Niech $\Gamma=(V,E)$ będzie skończonym, spójnym drzewem
prostym, a $e\colon V\to\mathbb Z$ funkcją wag.
Dla $v\in V$ bierzemy wiązkę $E(e_v)\to S^2$.
Dla każdej krawędzi $uv\in E$ wykonujemy dodatni plumbing
wiązek przypisanych do $u$ i $v$, używając rozłącznych
dysków bazowych przy różnych krawędziach. Po wygładzeniu
narożników otrzymujemy zwartą, zorientowaną, gładką
czterorozmaitość z brzegiem $P_\Gamma$.
\end{definicja}

W ogólnym grafie plumbingowym można także dopuszczać
krawędzie ze znakiem ujemnym oraz pętle; trzeba
wtedy ustalić dodatkowe dane orientacyjne i sposób
sklejenia. Znaki krawędzi nie są wagami Eulera
wierzchołków. W drzewie każdą konfigurację znaków
krawędzi można sprowadzić do samych znaków dodatnich
przez zmianę orientacji klas sfer: ustalamy znak
w korzeniu, a następnie wybieramy znak każdej
kolejnej sfery tak, aby jej przecięcie ze sferą
rodzica było dodatnie. Jedyna droga od korzenia
wyklucza konflikt wyborów. W tej części, jeżeli nie
zaznaczono inaczej, graf oznacza \emph{ważone drzewo} z dodatnim
lokalnym znakiem przecinania, tak jak w powyższej definicji.

\section{Forma przecięcia z grafu}
\label{sec:tg-plumbing-Q}

\begin{twierdzenie}[Forma przecięcia drzewa plumbingowego]
\label{thm:tg-plumbing-intersection}
Niech $\Gamma=(V,E)$ będzie skończonym, spójnym, ważonym
drzewem plumbingowym, a $P_\Gamma$ czterorozmaitością
z poprzedniej definicji. Wtedy $\pi_1(P_\Gamma)=1$ oraz
$H_2(P_\Gamma;\mathbb Z)\cong\mathbb Z^{|V|}$.
Ponadto $H_k(P_\Gamma;\mathbb Z)=0$ dla $k=1$
i wszystkich $k\ge3$.
W bazie $[S_v]$ wyznaczonej przez obrazy sekcji zerowych
macierz formy przecięcia wynosi
\[
Q_\Gamma=(q_{uv})_{u,v\in V},\qquad
q_{uv}=\begin{cases}
e_v,&u=v,\\
1,&uv\in E,\\
0,&u\ne v\text{ i }uv\notin E.
\end{cases}
\]
\end{twierdzenie}

\begin{proof}
Najpierw rozpatrujemy pojedynczy wierzchołek $v$.
Wiązka $E(e_v)$ ma retrakcję deformacyjną
na sekcję zerową $S^2$, daną przez mnożenie
wektorów we włóknach przez parametr z $[0,1]$.
Jest zatem jednospójna, a jej zredukowana
homologia jest równa $\mathbb Z$ w stopniu
dwa i zerowa w pozostałych stopniach.
Dla drzewa o co najmniej dwóch wierzchołkach
usuwamy liść $w$; pozostałe drzewo jest spójne.
Odtwarzając liść, dołączamy $E(e_w)$ do
$P_{\rm stary}$ wzdłuż dokładnie jednego bloku
$D^2\times D^2$; rozłączność dysków bazowych
wyklucza dalsze identyfikacje.
Przed wygładzeniem narożników części i ich
wspólny blok można traktować jako podkompleksy
wspólnej triangulacji zgodnej ze ścianami.
Ich otoczenia otwarte (względnie w całej
przestrzeni) mają te same typy homotopijne
co części, a przecięcie ma typ wspólnego bloku,
przez standardowe otoczenia podkompleksów
w podziale barycentrycznym. Wygładzenie
narożników nie zmienia typu topologicznego.
Możemy więc stosować twierdzenie van Kampena
i ciąg Mayera--Vietorisa. Przecięcie jest
ściągalne i spójne. Twierdzenie van Kampena
daje iloczyn wolny dwóch trywialnych grup.
Ciąg Mayera--Vietorisa w stopniu drugim ma postać
\[
0\longrightarrow H_2(P_{\rm stary};\mathbb Z)\oplus
H_2(E(e_w);\mathbb Z)\longrightarrow
H_2(P_{\rm nowy};\mathbb Z)\longrightarrow0.
\]
Indukcja względem liczby wierzchołków daje opis grupy
$H_2$ oraz bazę sekcji zerowych, ponieważ
izomorfizm jest indukowany przez włożenia obu
części. W pozostałych dodatnich stopniach
korzystamy z tej samej dokładności dla homologii
zredukowanej: zredukowana homologia przecięcia
jest zerowa w każdym stopniu, więc homologia
nowej części jest sumą prostą homologii obu
części. Otrzymujemy też żądane zanikanie.

Dla $v=u$ wiązka normalna sekcji $S_v$
pozostaje jej pierwotną wiązką wektorową:
w każdym wspólnym bloku lokalne współrzędne
normalne do $S_v$ są nadal jego współrzędnymi
włóknowymi, choć dla drugiej sfery są to
współrzędne bazowe. Funkcje przejścia wzdłuż
$S_v$ nie zmieniają się. Zatem samoprzecięcie
wynosi $e_v$ z rachunku klasy Eulera. Gdy $uv\in E$,
sekcje $S_u$ i $S_v$ przecinają się dokładnie w jednym
bloku plumbingowym i znak ich przecięcia, zgodnie z
przyjętą orientacją, jest dodatni. Gdy brak krawędzi,
sekcje można reprezentować rozłącznie. Otrzymujemy
wskazane elementy macierzy. Symetria macierzy wynika
z symetrii formy przecięcia w wymiarze cztery.
\end{proof}

\begin{przyklad}[Dwa wierzchołki]
Dla drzewa z jedną krawędzią i wagami $-2,-3$ dostajemy
\[
Q_\Gamma=\begin{pmatrix}-2&1\\1&-3\end{pmatrix},
\qquad \det Q_\Gamma=5.
\]
Dwie sfery bazowe mają samoprzecięcia $-2$ i $-3$ oraz
przecinają się dodatnio w jednym punkcie. Jak zobaczymy,
ze względu na niezerowy wyznacznik brzeg tej czterorozmaitości
jest wymierną sferą homologiczną i ma pierwszą grupę homologii
rzędu pięć.
\end{przyklad}

\begin{center}
\begin{tikzpicture}[scale=1,>=Stealth]
  \node[tgBlueVertex,minimum size=12mm] (a) at (-1.5,0) {$-2$};
  \node[tgRedVertex,minimum size=12mm] (b) at (1.5,0) {$-3$};
  \draw[tgTree] (a)--(b);
  \node[font=\footnotesize,text=tgNavy] at (0,-.95)
    {Jedna krawędź: liczba przecięcia $[S_a]\cdot[S_b]=1$.};
\end{tikzpicture}
\end{center}

\begin{uwaga}[Konstrukcja przez uchwyty]
Ten sam $P_\Gamma$ można opisać za pomocą ramowanego splotu
w $S^3$. Każdemu wierzchołkowi odpowiada niezawiązany okrąg
z ramowaniem $e_v$, a każdej krawędzi -- dodatnie splecenie
tych dwóch okręgów równe jeden. Dla drzewa można wybrać
standardowy splot plumbingowy, w którym pozostałe liczby
splecenia wynoszą zero. Dołączenie do $B^4$ uchwytów indeksu
dwa wzdłuż tego splotu daje dyfeomorficzną czterorozmaitość.
Identyfikację geometryczną przyjmujemy z klasycznej
konstrukcji diagramów Kirby'ego dla drzew plumbingowych:
Gompf--Stipsicz, \emph{4-Manifolds and Kirby Calculus},
sekcja 4.6, s.~129. Dla pojedynczego wierzchołka
rozkład bazy na dwa dyski daje kulę i jeden
uchwyt o ramowaniu $e_v$, tak jak w rozdziale
\ref{ch:kt-cztery}. Przy dołączaniu liścia
zamiana bazy i włókna daje nową krzywą
przyczepienia będącą małym meridianem
krzywej przyczepienia rodzica, z ramowaniem
liścia. Wybieramy te meridiany w rozłącznych
lokalnych otoczeniach. Nie utożsamiamy
dowolnych splotów o tej samej macierzy
splecenia: chodzi o konkretny standardowy
diagram z tej konstrukcji.
W szczególności macierz splecenia tego splotu jest macierzą
$Q_\Gamma$.
\end{uwaga}

\section{Homologia i forma splecenia brzegu}
\label{sec:tg-plumbing-boundary}

\begin{twierdzenie}[Homologia brzegu plumbingowania]
\label{thm:tg-plumbing-boundary}
Niech $\Gamma$ będzie skończonym, spójnym, ważonym drzewem,
a $Y_\Gamma:=\partial P_\Gamma$ brzegiem odpowiadającej
mu czterorozmaitości. Wtedy
\[
H_1(Y_\Gamma;\mathbb Z)\cong
\operatorname{coker}(Q_\Gamma:\mathbb Z^{|V|}\to\mathbb Z^{|V|}),
\qquad
H_2(Y_\Gamma;\mathbb Z)\cong\ker Q_\Gamma.
\]
Jeżeli $\det Q_\Gamma\ne0$, to $Y_\Gamma$ jest wymierną
sferą homologiczną oraz
$|H_1(Y_\Gamma;\mathbb Z)|=|\det Q_\Gamma|$.
Jeżeli $|\det Q_\Gamma|=1$, to $Y_\Gamma$ jest całkowitą
sferą homologiczną.
\end{twierdzenie}

\begin{proof}
Oznaczmy $W=P_\Gamma$, $Y=\partial W$ i $r=|V|$.
W poprzednim twierdzeniu wykazaliśmy, że
$H_1(W;\mathbb Z)=H_3(W;\mathbb Z)=0$
i $H_2(W;\mathbb Z)\cong\mathbb Z^r$.
O ile nie zaznaczono inaczej, współczynniki
homologii są w $\mathbb Z$.
Najpierw uzasadnijmy spójność brzegu.
Twierdzenie o uniwersalnych współczynnikach
daje $H^3(W)=0$ (zarówno $H_3(W)$,
jak i $\operatorname{Ext}(H_2(W),\mathbb Z)$
są zerowe) oraz $H^4(W)=0$.
Z dualności Poincarégo--Lefschetza
$H_1(W,Y)=H_0(W,Y)=0$.
Ciąg dokładny pary w stopniu zero daje więc
izomorfizm $H_0(Y)\to H_0(W)=\mathbb Z$.
W szczególności $Y$ jest niepusty i spójny.
W długim ciągu dokładnym pary $(W,Y)$ występuje fragment
\[
H_3(W,Y)\longrightarrow H_2(Y)\longrightarrow H_2(W)
\xrightarrow{\ j\ }H_2(W,Y)\longrightarrow H_1(Y)
\longrightarrow H_1(W).
\]
Z dualności Poincarégo--Lefschetza i uniwersalnego twierdzenia
o współczynnikach wynika
$H_3(W,Y)\cong H^1(W)=0$ oraz
$H_2(W,Y)\cong H^2(W)\cong\operatorname{Hom}(H_2(W),\mathbb Z)$.
W wybranych bazach odwzorowanie $j$ zapisuje się jako macierz
formy przecięcia $Q_\Gamma$: obraz klasy $x$ jest funkcjonałem
$y\mapsto x\cdot y$. Ciąg dokładny upraszcza się więc do
\[
0\longrightarrow H_2(Y)\longrightarrow\mathbb Z^r
\xrightarrow{Q_\Gamma}\mathbb Z^r
\longrightarrow H_1(Y)\longrightarrow0,
\]
co dowodzi dwóch pierwszych izomorfizmów. Jeśli $Q_\Gamma$
ma niezerowy wyznacznik, jego jądro nad $\mathbb Z$ jest zerowe,
a iloraz $\mathbb Z^r/Q_\Gamma\mathbb Z^r$ jest skończony
o liczbie elementów $|\det Q_\Gamma|$ (np. z postaci normalnej
Smitha). Dualność Poincarégo na zamkniętej, spójnej,
zorientowanej trójrozmaitości daje wtedy pozostałe grupy
homologii takie jak dla sfery nad $\mathbb Q$.
Przy wyznaczniku $\pm1$ również $H_1(Y)=0$, co daje
analogiczne stwierdzenie nad $\mathbb Z$.
\end{proof}

\begin{przyklad}[Macierz i relacje w pierwszej homologii]
Dla drzewa o wagach $-2,-3$ macierz z poprzedniego
przykładu daje prezentację
$H_1(Y;\mathbb Z)=\langle x,y\mid -2x+y=0,\ x-3y=0\rangle$.
Pierwsza relacja daje $y=2x$, a druga $-5x=0$.
Zatem $H_1(Y;\mathbb Z)\cong\mathbb Z/5\mathbb Z$.
Zauważmy, że wyznacznik macierzy wystarcza do obliczenia
\emph{rzędu} tej grupy, ale nie zawsze do wyznaczenia
jej struktury: przykładem będzie graf typu $D_4$.
\end{przyklad}

\begin{lemat}[Rozwinięcie wyznacznika wzdłuż liścia]
\label{lem:tg-plumbing-leaf-determinant}
Niech $\Gamma=(V,E)$ będzie skończonym ważonym drzewem,
$v\in V$ jego liściem, a $u$ jedynym sąsiadem $v$.
Załóżmy, że przecięcia przy krawędziach są równe $+1$
i oznaczmy wagę liścia przez $a=e_v$.
Wówczas
\[
\det Q_\Gamma
=a\det Q_{\Gamma\setminus\{v\}}
-\det Q_{\Gamma\setminus\{u,v\}},
\]
gdzie dla rozłącznej sumy drzew wyznacznik
jest iloczynem wyznaczników jej składowych,
a dla grafu pustego przyjmujemy wyznacznik równy $1$.
Wzór zachodzi także przy $a=0$ i wagach ujemnych.
\end{lemat}

\begin{proof}
Uporządkujmy wierzchołki tak, by $v$ był ostatni.
Wtedy jedynymi niezerowymi wyrazami
w ostatnim wierszu $Q_\Gamma$ są
$1$ w kolumnie $u$ i $a$ na przekątnej.
Rozwinięcie Laplace'a względem tego wiersza
daje składnik $a\det Q_{\Gamma-v}$ i składnik
odpowiadający krawędzi $uv$.
Niech $r=|V|$ i niech $u$ ma numer $j<r$.
Pierwsze rozwinięcie przy wyrazie w kolumnie $j$
daje znak $(-1)^{r+j}$. W otrzymanym minorze
kolumna odpowiadająca $v$ ma numer $r-1$,
a jej jedyny niezerowy wyraz równy $1$ leży
w wierszu $j$. Drugie rozwinięcie daje znak
$(-1)^{j+r-1}$. Iloczyn obu znaków jest równy
$(-1)^{2r+2j-1}=-1$, a pozostała macierz
jest macierzą $Q_{\Gamma-\{u,v\}}$.
Ten ostatni graf może być lasem,
a jego macierz jest wtedy blokowo diagonalna.
Nie dzieliliśmy przez $a$ ani przez inny wyznacznik,
więc tożsamość obejmuje przypadki osobliwe.
\end{proof}

\begin{przyklad}[Dwie niezależne metody dla $D_4$]
\label{ex:tg-plumbing-D4-leaf}
Dla grafu $D_4$ o centralnej wadze $-2$
i trzech liściach o wadze $-2$
usuwamy jeden liść $v$ i jego sąsiada,
czyli wierzchołek centralny $u$.
Pozostają odpowiednio łańcuch $A_3$
o wyznaczniku $-4$ i dwa odizolowane
wierzchołki $-2$ o wyznaczniku $4$.
Stąd lemat daje
$\det Q_{D_4}=(-2)(-4)-4=4$,
co zgadza się z rachunkiem przez dopełnienie Schura
w sekcji \ref{sec:tg-plumbing-Dn}.
Porównanie obu metod jest szczególnie
przydatne przy audycie dużych drzew,
w których błędny znak zmieniłby
homologię brzegu.
\end{przyklad}

\begin{stwierdzenie}[Forma splecenia brzegu]
Niech $\Gamma$ będzie jak wyżej i niech $Q_\Gamma$ będzie
nieosobliwa nad $\mathbb Q$. Dla $Y=\partial P_\Gamma$
forma splecenia na skończonej grupie
$H_1(Y;\mathbb Z)\cong\mathbb Z^r/Q_\Gamma\mathbb Z^r$
ma, przy konwencji orientacyjnej brzegu jako $\partial P_\Gamma$,
postać
\[
\lambda_Y([x],[y])\equiv-\,x^TQ_\Gamma^{-1}y
\pmod{\mathbb Z}, \qquad x,y\in\mathbb Z^r.
\]
\end{stwierdzenie}

\begin{proof}
Macierz $Q_\Gamma$ jest całkowita, symetryczna i nieosobliwa.
Jeżeli zastąpimy $x$ przez $x+Q_\Gamma z$, wyrażenie po prawej
zmienia się o $-z^Ty\in\mathbb Z$; analogicznie dla $y$.
Forma jest zatem dobrze określona modulo $\mathbb Z$.
Jej niezdegenerowanie wynika z tego, że jeśli
$x^TQ_\Gamma^{-1}y\in\mathbb Z$ dla każdego $y\in\mathbb Z^r$,
to $Q_\Gamma^{-1}x\in\mathbb Z^r$, a więc
$[x]=0$ w grupie ilorazowej. Identyfikację z geometryczną formą splecenia
otrzymujemy bez ponownego skrótowego rachunku
znaku: stosujemy twierdzenie o formie splecenia
z macierzy chirurgii z sekcji
\ref{sec:kt-linking-pairing} do standardowego
diagramu Kirby'ego drzewa. Jego macierz
splecenia jest właśnie $Q_\Gamma$.
Identyfikacja generatorów w prezentacji jest
zgodna z poprzednim dowodem: meridiany są
brzegami współrdzeni uchwytów, a współrdzenie
tworzą bazę $H_2(W,Y)$ dualną względem
przecięcia do domkniętych rdzeni (sekcji
zerowych). Dowód w przytoczonej sekcji,
przez homomorfizm Bocksteina i rachunek
iloczynu kohomologicznego na klasie
fundamentalnej pary $(W,Y)$, daje dokładnie
znak minus dla orientacji brzegowej.
Nie wystarcza tu sama zgodność macierzy
czy grupy homologii.
\end{proof}

\begin{uwaga}
Przy odwróceniu orientacji $Y$ forma splecenia zmienia znak.
Nie należy więc porównywać formuł $Q^{-1}$ i $-Q^{-1}$
bez sprawdzenia orientacji i definicji dodatniego meridianu.
Sama grupa $H_1(Y;\mathbb Z)$ na ogół nie rozróżnia
różnych trójrozmaitości.
\end{uwaga}

\section{Łańcuchy ważone, ułamki łańcuchowe i przestrzenie soczewkowe}
\label{sec:tg-plumbing-chains}

Najłatwiejsze rachunki wykonujemy dla grafów liniowych.
Złożenie kolejnych bloków można wtedy opisać zarówno
macierzą trójdiagonalną, jak i ułamkiem łańcuchowym.

\begin{definicja}[Łańcuch plumbingowy]
Niech $r\ge1$ i niech $a_1,\dots,a_r\in\mathbb Z$.
Przez $C(a_1,\dots,a_r)$ oznaczamy ważone drzewo o
wierzchołkach $v_1,\dots,v_r$, krawędziach
$v_iv_{i+1}$ dla $1\le i<r$ i wagach $e(v_i)=a_i$.
Jego macierz przecięcia oznaczamy przez $Q(a_1,\dots,a_r)$.
\end{definicja}

\begin{lemat}[Rekurencja wyznaczników]
\label{lem:tg-plumbing-chain-det}
Dla $r\ge2$ wyznaczniki
$d_r:=\det Q(a_1,\dots,a_r)$ spełniają
\[
d_r=a_r d_{r-1}-d_{r-2},\qquad
 d_0=1,\quad d_1=a_1.
\]
\end{lemat}
\begin{proof}
Ostatni wiersz macierzy $Q(a_1,\dots,a_r)$ ma
jedynie dwa potencjalnie niezerowe wyrazy: $1$ na pozycji
$r-1$ oraz $a_r$ na pozycji $r$. Rozwijamy wyznacznik
Laplace'a wzdłuż ostatniego wiersza. Składnik z $a_r$
równa się $a_r d_{r-1}$. Minor odpowiadający wyrazowi $1$
na pozycji $r-1$ ma w ostatniej kolumnie tylko jedną
jedynkę; po drugim rozwinięciu dostajemy $-d_{r-2}$.
Dla $r=2$ otrzymujemy bezpośrednio
$d_2=a_1a_2-1$, co zgadza się z rekurencją.
\end{proof}

\begin{stwierdzenie}[Ułamek łańcuchowy dla chirurgii]
\label{st:tg-plumbing-chain-surgery}
Niech $a_1,\dots,a_r\in\mathbb Z$ będą liczbami,
dla których wszystkie mianowniki w poniższym wyrażeniu
są niezerowe. Niech $L$ będzie standardowym ramowanym
splotem będącym łańcuchem niezawiązanych okręgów o
ramowaniach $a_1,\dots,a_r$. Operacje \emph{slam-dunk}
(racjonalnej eliminacji kolejnego torusa pełnego)
prowadzą do wymiernej chirurgii na jednym węźle trywialnym
o współczynniku
\[
[a_1,\dots,a_r]^-:=a_1-\cfrac1{a_2-
\cfrac1{\ddots-\cfrac1{a_r}}}.
\]
W szczególności brzeg odpowiadającego plumbingowania
jest przestrzenią soczewkową, gdy licznik tego ułamka
jest niezerowy (z uwzględnieniem orientacji i konwencji
parametrów przestrzeni soczewkowych).
\end{stwierdzenie}

\begin{proof}[Rachunek współczynnika]
Stosujemy geometryczny ruch slam-dunk w wersji
dla całkowitego współczynnika $a$ na składowej
i wymiernego niezerowego współczynnika $b$ na
jej lokalnym meridianie. Ruch zachowuje
orientowany wynik chirurgii i zastępuje oba
współczynniki jednym $a-1/b$.
Jest to zewnętrzny wynik rachunku chirurgii:
R.~E. Gompf, A.~I. Stipsicz,
\emph{4-Manifolds and Kirby Calculus}, AMS, 1999,
s.~163--164 (slam-dunk i ćwiczenie 5.3.9).
Nie wyprowadzamy dyfeomorfizmu z eliminacji
macierzowej.

W standardowym diagramie łańcucha ostatnia
składowa jest lokalnym meridianem przedostatniej,
a jej usunięcie nie zmienia pozostałych składowych.
Po pierwszym ruchu przedostatnia składowa ma
współczynnik $a_{r-1}-1/a_r$. Jest ona nadal
lokalnym meridianem poprzedniej składowej,
której współczynnik pozostaje całkowity.
Indukcyjnie można więc powtarzać ruch,
zawsze z niezerowym mianownikiem na mocy założenia.
Pozostaje węzeł trywialny z podanym
współczynnikiem. Wypełnienie jego zewnętrza,
które jest torusem pełnym, stanowi sklejenie
dwóch torusów pełnych. Przy niezerowym liczniku
jest to przestrzeń soczewkowa; przy liczniku
zerowym otrzymujemy $S^1\times S^2$.
Nie utożsamiamy przestrzeni soczewkowych
na podstawie samej homologii.
\end{proof}

\begin{stwierdzenie}[Homologia i ułamek łańcuchowy]
\label{st:tg-plumbing-chain-homology}
Dla dowolnych całkowitych wag $a_1,\ldots,a_r$
oznaczmy $D_i=\det Q(a_i,\ldots,a_r)$,
$D_{r+1}=1$, $D_{r+2}=0$. Wtedy
\[
D_i=a_iD_{i+1}-D_{i+2},\qquad
\gcd(D_i,D_{i+1})=1.
\]
Pierwsza homologia brzegu jest cykliczna:
\[
H_1(Y;\mathbb Z)\cong\mathbb Z/(D_1).
\]
Dla $D_1=0$ oznacza to grupę $\mathbb Z$.
Jeżeli ułamek z poprzedniego stwierdzenia
jest określony, wynosi $D_1/D_2$.
Dla $D_1\ne0$ na generatorze $x$ odpowiadającym
pierwszemu wierzchołkowi zachodzi
$\lambda_Y(x,x)\equiv-D_2/D_1\pmod{\mathbb Z}$.
\end{stwierdzenie}
\begin{proof}
Rekurencja wynika z lematu
\ref{lem:tg-plumbing-chain-det}, po odwróceniu
kolejności wierzchołków. Z własności
największego wspólnego dzielnika
$\gcd(D_i,D_{i+1})=\gcd(D_{i+1},D_{i+2})$
oraz warunków końcowych wynika względna pierwszość.
Jeżeli kolejne mianowniki są niezerowe,
indukcja od końca daje
$[a_i,\ldots,a_r]^-=D_i/D_{i+1}$.

Dla $r=1$ prezentacja homologii jest bezpośrednia.
Dla $r\ge2$ oznaczmy generatory przez $x_j$,
a formalnie $x_0=x_{r+1}=0$.
Relacje $x_{j-1}+a_jx_j+x_{j+1}=0$
pozwalają wyeliminować kolejno
$x_2,\ldots,x_r$ ze współczynnikiem jednostkowym.
Pisząc $d_k=\det Q(a_1,\ldots,a_k)$
i $d_0=1$, otrzymujemy
$x_j=(-1)^{j-1}d_{j-1}x_1$ przez indukcję
z rekurencji wyznaczników.
Ostatnia relacja jest równoważna $d_rx_1=0$;
nie pozostają żadne inne relacje,
gdyż każda eliminacja była odwracalna
w prezentacji grupy abelowej.
Stąd $H_1(Y;\mathbb Z)\cong\mathbb Z/(D_1)$.
Jeśli $D_1\ne0$, element $(1,1)$ macierzy
odwrotnej jest równy $D_2/D_1$ na mocy
wzoru przez dopełnienia algebraiczne.
Stosujemy wzór formy splecenia z sekcji
\ref{sec:tg-plumbing-boundary}.
W szczególności względna pierwszość
$D_1,D_2$ zapewnia niezdegenerowanie
formy na tej grupie cyklicznej.
\end{proof}

\begin{przyklad}[Osobliwy łańcuch]
Dla $C(1,1)$ wyznacznik wynosi zero,
a ułamek łańcuchowy jest równy $1-1=0$.
Zatem brzeg jest dyfeomorficzny z
$S^1\times S^2$, a pierwsza homologia
jest izomorficzna z $\mathbb Z$.
Natomiast dla $C(0,0)$ wyznacznik wynosi
$-1$, chociaż zwykły ułamek łańcuchowy
nie jest określony. Pierwsza homologia
jest zerowa, a brzeg jest $S^3$
z wcześniejszego przykładu chirurgii
na splocie Hopfa w rozdziale
\ref{ch:kt-cztery}. Pokazuje to, że
rachunek homologii nie wymaga niezerowych
mianowników pośrednich.
\end{przyklad}

\begin{uwaga}[Orientacja współczynnika]
Dla łańcucha ujemnych wag $-b_i$, gdzie $b_i\ge2$,
wartość $[-b_1,\dots,-b_r]^-$ jest ujemna,
a wszystkie mianowniki pośrednie są niezerowe.
Istotnie, niech $t_r=-b_r$ oraz
$t_i=-b_i-1/t_{i+1}$ dla $i=r-1,\dots,1$.
Mamy $t_r\le -2<-1$. Jeżeli $t_{i+1}<-1$,
to $0<-1/t_{i+1}<1$, skąd
$t_i<-b_i+1\le-1$. Indukcja od końca
daje $t_i<-1$ dla każdego $i$ i uzasadnia
stosowalność stwierdzenia
\ref{st:tg-plumbing-chain-surgery}.
Przy przejściu od tej liczby do oznaczenia $L(p,q)$
trzeba ustalić znak licznika oraz orientację brzegu;
różne podręczniki stosują konwencje $L(p,q)$ i
$L(p,q^{-1})$ lub ich odbicia lustrzane.
Poniższe obliczenia homologii nie zależą od tych wyborów.
\end{uwaga}

\begin{twierdzenie}[Łańcuch typu $A_n$]
\label{thm:tg-plumbing-An}
Niech $n\ge1$, a $\Gamma=A_n$ będzie łańcuchem
$n$ wierzchołków o wagach $-2$. Niech $P_{A_n}$
będzie jego czterorozmaitością plumbingową. Wtedy
\[
\det Q_{A_n}=(-1)^n(n+1),\qquad
H_1(\partial P_{A_n};\mathbb Z)\cong\mathbb Z/(n+1)\mathbb Z.
\]
Forma przecięcia $Q_{A_n}$ jest ujemnie określona.
Brzeg $P_{A_n}$ jest przestrzenią soczewkową o pierwszej
grupie homologii rzędu $n+1$.
\end{twierdzenie}
\begin{proof}
Niech $d_n$ będzie wyznacznikiem odpowiedniej macierzy.
Z lematu o wyznaczniku łańcucha mamy
$d_n=-2d_{n-1}-d_{n-2}$, gdzie $d_0=1$, $d_1=-2$.
Podstawienie $d_n=(-1)^n(n+1)$ spełnia rekurencję
i warunki początkowe, zatem wzór jest prawdziwy dla
wszystkich $n$.

Aby wyznaczyć grupę, oznaczmy jej generatory odpowiadające
kolejnym wierzchołkom przez $x_1,\dots,x_n$. Dla $n=1$ grupa jest bezpośrednio
$\mathbb Z/2\mathbb Z$, więc dalej przyjmijmy $n\ge2$.
Relacje
macierzowe mają postać
$-2x_1+x_2=0$ na początku,
$x_{i-1}-2x_i+x_{i+1}=0$ dla $1<i<n$ oraz
$x_{n-1}-2x_n=0$ na końcu.
Przyjmując formalnie $x_0=x_{n+1}=0$, zapisujemy je
jednolicie jako $x_{i-1}-2x_i+x_{i+1}=0$.
Indukcyjnie otrzymujemy $x_i=i x_1$, a ostatnia relacja
daje $(n+1)x_1=0$. Grupa jest więc cykliczna i ma
rząd $n+1$ na mocy poprzedniego twierdzenia.

Wszystkie wiodące minory główne macierzy $Q_{A_n}$
mają wyznaczniki $(-1)^k(k+1)$. Z kryterium Sylvestera
wynika, że $Q_{A_n}$ jest ujemnie określona.
Wreszcie współczynnik chirurgii wyrażony przez ułamek
łańcuchowy wynosi
$[-2,\dots,-2]^-=-(n+1)/n$ dla $n\ge1$:
jest to iloraz wyznaczników $d_n/d_{n-1}$
ze stwierdzenia \ref{st:tg-plumbing-chain-homology},
a wszystkie wyznaczniki ogonów są niezerowe.
Dlatego brzeg otrzymuje się przez wymierną chirurgię
na węźle trywialnym, co kończy dowód.
\end{proof}

\begin{center}
\begin{tikzpicture}[scale=.95,line cap=round]
  \foreach \x in {0,1,2,3,4}{
    \node[tgBlueVertex,minimum size=9mm] (v\x) at ({1.45*\x},0) {$-2$};
  }
  \foreach \x in {0,1,2,3}{
    \pgfmathtruncatemacro{\y}{\x+1}
    \draw[tgTree] (v\x)--(v\y);
  }
  \node[text=tgNavy,font=\small] at (2.9,-.85) {Przykład: $A_5$, wyznacznik $-6$};
\end{tikzpicture}
\end{center}

\begin{stwierdzenie}[Forma splecenia brzegu łańcucha typu $A_n$]
\label{st:tg-plumbing-An-linking}
Niech $n\ge1$, a $P_{A_n}$ będzie ujemnym
plumbingowaniem łańcucha $A_n$ z wagami $-2$.
Przy orientacji $Y=\partial P_{A_n}$
i generatorze $x=[x_1]$ grupy
$H_1(Y;\mathbb Z)\cong\mathbb Z/(n+1)\mathbb Z$,
związanym z pierwszym wierzchołkiem,
forma splecenia torsyjnego spełnia
\[
\lambda_Y(x,x)\equiv\frac{n}{n+1}\pmod{\mathbb Z}.
\]
W szczególności dla $A_1,A_2,A_3$
otrzymujemy odpowiednio wartości
$1/2$, $2/3$ oraz $3/4$ modulo $\mathbb Z$.
\end{stwierdzenie}
\begin{proof}
W poprzednim twierdzeniu~\ref{thm:tg-plumbing-An}
wykazaliśmy, że $x_1$ generuje
grupę cykliczną rzędu $n+1$.
Dla macierzy przecięcia $Q_{A_n}$
wzór na element macierzy odwrotnej
przez dopełnienie algebraiczne daje
\[
(Q_{A_n}^{-1})_{11}
=\frac{\det Q_{A_{n-1}}}{\det Q_{A_n}}
=\frac{(-1)^{n-1}n}{(-1)^n(n+1)}
=-\frac{n}{n+1},
\]
gdzie przy $n=1$ używamy
konwencji $\det Q_{A_0}=1$.
Z twierdzenia o formie splecenia brzegu
plumbingowania wynika
$\lambda_Y(x,x)=-(Q_{A_n}^{-1})_{11}\bmod\mathbb Z$,
co daje tezę. Wartość ma dokładnie
mianownik $n+1$, więc ta forma jest
niezdegenerowana na grupie cyklicznej
zgodnie z ogólnym twierdzeniem.
\end{proof}

\section{Rozgałęzienia: typ \texorpdfstring{$D_n$}{Dn} i różne grupy brzegowe}
\label{sec:tg-plumbing-Dn}

\begin{przyklad}[Graf $D_4$: wyznacznik nie wystarcza]
Rozważmy gwiazdę z czterema wierzchołkami, w której
wierzchołek centralny łączy się z trzema liśćmi,
a wszystkie wagi wynoszą $-2$. W kolejności
(centrum, trzy liście) forma przecięcia ma macierz
\[
Q_{D_4}=\begin{pmatrix}
-2&1&1&1\\1&-2&0&0\\1&0&-2&0\\1&0&0&-2
\end{pmatrix}.
\]
Blok odpowiadający trzem liściom jest równy $-2I_3$.
Jego dopełnienie Schura to
$-2-(1,1,1)(-\tfrac12I_3)(1,1,1)^T=-\tfrac12$.
Dlatego $\det Q_{D_4}=(-2)^3(-\tfrac12)=4$.

Rachunek modulo dwa pokazuje jednak, że
$H_1(\partial P_{D_4};\mathbb Z)$ nie jest grupą
$\mathbb Z/4\mathbb Z$. Rzeczywiście, macierz
$Q_{D_4}$ po redukcji modulo dwa ma rząd $2$, więc
\[
\dim_{\mathbb F_2}
(H_1(\partial P_{D_4};\mathbb Z)\otimes\mathbb F_2)=2.
\]
Grupa rzędu cztery z dwoma niezależnymi elementami
modulo dwa musi być izomorficzna z
$\mathbb Z/2\mathbb Z\oplus\mathbb Z/2\mathbb Z$.
Pokazuje to, że rząd $|\det Q|$ nie określa całej
grupy; potrzebna jest np. postać normalna Smitha.
\end{przyklad}

\begin{center}
\begin{tikzpicture}[scale=.9]
  \node[draw=tgPurple,fill=tgPurple!14,circle,
    minimum size=10mm] (c) at (0,0) {$-2$};
  \node[tgBlueVertex,minimum size=10mm] (a) at (-1.6,0) {$-2$};
  \node[tgRedVertex,minimum size=10mm] (b) at (1.6,0) {$-2$};
  \node[tgTealVertex,minimum size=10mm] (d) at (0,1.35) {$-2$};
  \draw[tgEdge] (a)--(c)--(b) (c)--(d);
  \node[font=\footnotesize,text=tgNavy] at (0,-.9)
    {$D_4$: centralny wierzchołek o stopniu trzy};
\end{tikzpicture}
\end{center}

\begin{twierdzenie}[Cała rodzina $D_n$]
\label{thm:tg-plumbing-Dn}
Niech $n\ge4$, a wszystkie wagi standardowego
grafu Dynkina $D_n$ będą równe $-2$.
Forma $Q_{D_n}$ jest ujemnie określona i
$\det Q_{D_n}=(-1)^n4$. Pierwsza grupa homologii brzegu jest
izomorficzna z $\mathbb Z/4\mathbb Z$
dla nieparzystego $n$, a z
$(\mathbb Z/2\mathbb Z)^2$ dla parzystego $n$.
\end{twierdzenie}
\begin{proof}
Po usunięciu centralnego wierzchołka
pozostają ramiona długości $n-3,1,1$.
Oznaczmy ich blok diagonalny przez $B$.
Każdy blok jest ujemnie określony na mocy
twierdzenia \ref{thm:tg-plumbing-An}.
Wektor $c$ łączący centrum z ramionami
ma jedynkę przy pierwszym wierzchołku
każdego ramienia. Z rachunku odwrotności
macierzy $A_m$ w stwierdzeniu
\ref{st:tg-plumbing-An-linking} otrzymujemy
\[
-2-c^TB^{-1}c
=-2+\frac{n-3}{n-2}+\frac12+\frac12
=-\frac1{n-2}.
\]
Ustawiając centrum jako ostatni wierzchołek,
zapisujemy $Q_{D_n}=\begin{psmallmatrix}B&c\\c^T&-2\end{psmallmatrix}$.
Dla $U=\begin{psmallmatrix}I&-B^{-1}c\\0&1\end{psmallmatrix}$
bezpośrednie mnożenie daje
$U^T Q_{D_n} U=B\oplus[-1/(n-2)]$.
Macierz $U$ jest wymierna i ma wyznacznik jeden.
Oba składniki są ujemnie określone.
Ponieważ $\det B=(-1)^{n-3}(n-2)\cdot4$,
wyznacznik całej macierzy wynosi $(-1)^n4$.
Aby rozpoznać strukturę grupy,
przechodzimy do prezentacji jej generatorów: wzdłuż długiego
ramienia relacje $-2x_i+x_{i-1}+x_{i+1}=0$ wyrażają
kolejne $x_i$ przez generator $x$ przy wolnym końcu.
Numerujemy wierzchołki od wolnego końca
do centrum, tak że $x_1=x$ i $x_{n-2}$ jest
generatorem centralnym. Przyjmujemy $x_0=0$.
Z rekurencji $x_{i+1}=2x_i-x_{i-1}$
dla $1\le i\le n-3$ wynika indukcyjnie
$x_i=ix$ dla $1\le i\le n-2$.
Jeśli $u,v$ oznaczają dwa krótkie liście, relacje przy
nich dają $2u=2v=(n-2)x$, zaś relacja centralna daje
$u+v=(n-1)x$. Eliminacja $v$ prowadzi zatem do
\[
H_1(\partial P_{D_n};\mathbb Z)
\cong \langle x,u\mid 2u=(n-2)x,\ 2x=0\rangle.
\]
Dla $n$ parzystego pierwsza relacja sprowadza się do
$2u=0$, więc otrzymujemy $(\mathbb Z/2)^2$.
Dla $n$ nieparzystego jest równoważna $2u=x$,
a wtedy $4u=0$ i otrzymujemy $\mathbb Z/4$.
Eliminacje są odwracalne w prezentacji:
$2v=2(n-1)x-2u=nx$ daje wraz z
$2v=(n-2)x$ dokładnie relację $2x=0$.
Nie pozostają żadne dalsze relacje.
\end{proof}

\begin{stwierdzenie}[Forma splecenia dla $D_n$]
\label{st:tg-plumbing-Dn-linking}
W oznaczeniach dowodu twierdzenia
\ref{thm:tg-plumbing-Dn}, gdzie $x$ odpowiada
wolnemu końcowi długiego ramienia, a $u$
jednemu z krótkich liści, forma splecenia
na brzegu ma wartości
\[
\lambda(x,x)=0,\qquad
\lambda(x,u)=\tfrac12,\qquad
\lambda(u,u)=\tfrac n4\pmod{\mathbb Z}.
\]
\end{stwierdzenie}
\begin{proof}
Niech $m=n-3$. Dla dodatniej macierzy
$C_m=-Q_{A_m}$ rozwiązanie $C_m z=e_1$
ma współrzędne $z_i=(m+1-i)/(m+1)$:
wzór obowiązuje dla $1\le i\le m$;
osobno przyjmujemy $z_0=z_{m+1}=0$.
Wówczas $2z_1-z_2=1$ (także dla $m=1$,
gdy $z_2=0$), a dla $2\le i\le m$ mamy
$2z_i-z_{i-1}-z_{i+1}=0$.
Nieosobliwość $C_m$ zapewnia jednoznaczność
rozwiązania.
W szczególności $(C_m^{-1})_{m1}=1/(m+1)$,
a przez symetrię łańcucha
$(C_m^{-1})_{mm}=m/(m+1)$.
Dla $-Q_{D_n}$ dopełnienie Schura centrum
wynosi $1/(n-2)$. Niech $T=-B$ i niech $c$ będzie wektorem
połączeń z dowodu poprzedniego twierdzenia.
Po ustawieniu centrum na końcu
$-Q_{D_n}=\begin{psmallmatrix}T&-c\\-c^T&2\end{psmallmatrix}$.
Aby obliczyć blok odwrotności na ramionach,
rozwiązujemy układ $Tz-ct=y$ i $-c^Tz+2t=0$
dla dowolnego wektora $y$ z przestrzeni ramion.
Z pierwszej równości $z=T^{-1}y+T^{-1}ct$.
Druga daje
$(2-c^TT^{-1}c)t=c^TT^{-1}y$.
Ponieważ dopełnienie Schura wynosi
$1/(n-2)$, podstawienie wyznacza szukany blok:
\[
T^{-1}+(n-2)T^{-1}cc^TT^{-1}.
\] Zatem wpisy odpowiadające
parom $(x,x),(x,u),(u,u)$ wynoszą
odpowiednio
\[
\frac{n-3}{n-2}+\frac1{n-2}=1,\quad
\frac12,\quad
\frac12+\frac{n-2}4=\frac n4.
\]
Redukując modulo $\mathbb Z$ i stosując
wzór formy splecenia z sekcji
\ref{sec:tg-plumbing-boundary}, otrzymujemy
tezę. Dla $n$ nieparzystego $x=2u$,
a $u$ ma rząd cztery. Rzeczywiście,
z dwuliniowości otrzymujemy
$\lambda(2u,u)=n/2\equiv1/2\pmod{\mathbb Z}$
oraz $\lambda(2u,2u)=n\equiv0\pmod{\mathbb Z}$,
zgodnie z obliczonymi wpisami.
Jeśli $\lambda(ku,u)=kn/4\equiv0\pmod{\mathbb Z}$,
to $4\mid kn$, a ponieważ $n$ jest nieparzyste,
$4\mid k$. Zatem $ku=0$ i forma jest
niezdegenerowana również w przypadku nieparzystym. Dla $n$ parzystego
macierz na generatorach $x,u$ ma wpisy
w $\frac12\mathbb Z/\mathbb Z$.
Stosujemy izomorfizm grup addytywnych
$\frac12\mathbb Z/\mathbb Z\to\mathbb F_2$,
$[k/2]\mapsto k\bmod 2$.
Nie jest to mnożenie przez dwa w grupie
$\mathbb Q/\mathbb Z$, które tutaj dawałoby
zero. Otrzymujemy macierz
$\begin{psmallmatrix}0&1\\1&n/2\end{psmallmatrix}$
o wyznaczniku $1$ w $\mathbb F_2$.
Dla $n\equiv0\pmod4$ jest to forma alternująca:
każdy element ma zerowe samosplecenie.
Dla $n\equiv2\pmod4$ mamy natomiast
$\lambda(u,u)=1/2$ i
$\lambda(x+u,x+u)=1/2$ modulo $\mathbb Z$,
a $\lambda(u,x+u)=1/2+n/4=0$ modulo $\mathbb Z$.
W bazie $u,x+u$ forma jest więc diagonalna
z obydwoma wpisami $1/2$.
Te dwie klasy izometrii są różne, gdyż
warunek zanikania wszystkich samospleceń
jest zachowywany przez izometrię.
\end{proof}

\section{Graf \texorpdfstring{$E_8$}{E8} i jego macierz przecięcia}
\label{sec:tg-plumbing-E8}

Najważniejszym przykładem jest rozgałęzione drzewo z ośmioma
wierzchołkami, występujące również jako diagram Dynkina
układu pierwiastków $E_8$. Wagi wszystkich wierzchołków
przyjmujemy równe $-2$.

\begin{definicja}[Ujemna forma $E_8$]
Niech $\Gamma_{E_8}$ będzie drzewem o wierzchołkach
$v_1,\ldots,v_8$, którego krawędzie to
$v_iv_{i+1}$ dla $1\le i\le6$ oraz $v_3v_8$.
Niech $Q_{E_8}$ ma elementy diagonalne równe $-2$,
pozadiagonalne równe $1$ dla sąsiednich wierzchołków
i zero w pozostałych miejscach. \emph{Ujemną formą
$E_8$} nazywamy formę całkowitą o macierzy $Q_{E_8}$.
Plumbing odpowiadający temu grafowi oznaczamy
przez $P_{E_8}$.
\end{definicja}

\begin{center}
\begin{tikzpicture}[scale=1.0,line cap=round]
  \foreach \i in {1,...,7}{
    \node[tgBlueVertex,minimum size=10mm] (v\i) at ({1.15*(\i-1)},0) {$-2$};
  }
  \foreach \i in {1,...,6}{
    \pgfmathtruncatemacro{\j}{\i+1}
    \draw[tgTree] (v\i)--(v\j);
  }
  \node[tgRedVertex,minimum size=10mm] (v8) at (2.30,1.40) {$-2$};
  \draw[tgAccent] (v3)--(v8);
  \node[font=\footnotesize,text=tgNavy] at (3.45,-.85)
    {Drzewo $E_8$: długości ramion $2$, $4$ i $1$.};
\end{tikzpicture}
\end{center}

\begin{twierdzenie}[Rachunek macierzy $E_8$]
\label{thm:tg-E8-unimodular}
Forma $Q_{E_8}$ jest parzysta, unimodularna i ujemnie
określona. W szczególności
\[
\det Q_{E_8}=1,\qquad \operatorname{sign}(Q_{E_8})=-8,
\qquad H_1(\partial P_{E_8};\mathbb Z)=0.
\]
\end{twierdzenie}

\begin{proof}
\emph{Krok pierwszy: parzystość.}
Dla $x=(x_v)_{v\in V}\in\mathbb Z^8$ mamy
\[
x^TQ_{E_8}x=-2\sum_vx_v^2+2\sum_{uv\in E}x_ux_v.
\]
Liczba ta jest parzysta, ponieważ oba składniki są
podzielne przez dwa. Jest to definicja parzystości
całkowitej formy biliniowej symetrycznej.

\emph{Krok drugi: wyznacznik i określoność.}
Usuńmy z drzewa wierzchołek trójwalentny $v_3$.
Powstają trzy rozłączne łańcuchy o długościach odpowiednio
$2$, $4$ i $1$, z macierzami $Q_{A_2},Q_{A_4},Q_{A_1}$
z twierdzenia \ref{thm:tg-plumbing-An}.
Każda z tych macierzy jest ujemnie określona.
Ich wyznaczniki wynoszą odpowiednio $3$, $5$ i $-2$.
Wszystkie trzy bloki razem mają więc wyznacznik $-30$.

Obliczmy dopełnienie Schura odpowiadające środkowemu
wierzchołkowi. Dla ujemnego łańcucha $A_m$ z
wpisami $-2$ na przekątnej i $1$ na sąsiednich
pozycjach wartość $(A_m^{-1})_{11}$ przy jego końcu
wynosi $-m/(m+1)$. Aby to wykazać, wystarczy użyć
wzoru na element macierzy odwrotnej jako iloraz
minorów: po usunięciu pierwszego wiersza i kolumny
pozostaje $A_{m-1}$, zatem
$(A_m^{-1})_{11}=\det A_{m-1}/\det A_m=-m/(m+1)$.
Wynik pozostaje taki sam, gdy ramię zapisujemy w
odwrotnej kolejności.

Dopełnienie Schura macierzy z trzema ramionami wynosi
\[
S=-2-\left(-\frac23-\frac45-\frac12\right)
=-2+\frac23+\frac45+\frac12=-\frac1{30}.
\]
Dokładniej, po ustawieniu centrum jako ostatniego
wierzchołka zapisujemy macierz w postaci
$\begin{psmallmatrix}B&c\\c^T&-2\end{psmallmatrix}$,
gdzie $B$ jest sumą bloków ramion.
Macierz
$U=\begin{psmallmatrix}I&-B^{-1}c\\0&1\end{psmallmatrix}$
ma wyznacznik jeden i spełnia
\[
U^TQ_{E_8}U=B\oplus[S].
\]
W tym wzorze $Q_{E_8}$ zapisujemy w przestawionej
kolejności wierzchołków. Stąd
$\det Q_{E_8}=(-30)(-1/30)=1$.
Z prawa bezwładności Sylvestera, ponieważ
wszystkie bloki ramion i $S$ są ujemnie
określone, wynika ujemna określoność
całej macierzy $Q_{E_8}$.
Kongruencja $U$ jest wymierna, nie całkowita;
nie wyciągamy z niej wniosku o rozkładzie kraty
na całkowitą sumę prostą. Macierz $Q_{E_8}$ ma zatem osiem
ujemnych wartości własnych i żadnej dodatniej,
więc jej sygnatura jest równa $-8$.

\emph{Krok trzeci: brzeg.}
Wyznacznik jest równy $1$, więc macierz $Q_{E_8}$
należy do $GL_8(\mathbb Z)$. Twierdzenie
\ref{thm:tg-plumbing-boundary} daje
$H_1(\partial P_{E_8};\mathbb Z)=0$.
Razem z dualnością Poincarégo w wymiarze trzy
pokazuje to, że brzeg jest całkowitą sferą homologiczną.
\end{proof}

\begin{uwaga}[Niezależny rachunek przez minory główne]
W pierwotnej kolejności $v_1,\ldots,v_8$ pierwsze
siedem wiodących minorów głównych pochodzi
z łańcuchów $A_k$ i wynosi $(-1)^k(k+1)$
dla $1\le k\le7$.
Stosując lemat \ref{lem:tg-plumbing-leaf-determinant}
do liścia $v_8$ i sąsiada $v_3$, otrzymujemy
\[
\det Q_{E_8}=(-2)\det Q_{A_7}
-\det Q_{A_2}\det Q_{A_4}=16-15=1.
\]
Znaki wszystkich ośmiu minorów głównych
są więc naprzemienne, zaczynając od znaku
ujemnego. Kryterium Sylvestera potwierdza
ujemną określoność niezależnie od rachunku
dopełnienia Schura.
\end{uwaga}

\begin{uwaga}[Dlaczego to wyjątkowy przykład]
Wśród dodatnio określonych parzystych krat unimodularnych
najmniejsza możliwa dodatnia ranga wynosi osiem.
Istotnie, klasyczny lemat van der Blija mówi,
że sygnatura parzystej formy unimodularnej
jest podzielna przez osiem; dla formy dodatnio
określonej jest ona równa randze.
Przyjmujemy ten wynik z literatury:
A.~Ranicki, J.~Roe, \emph{Surgery for Amateurs},
wersja z 19 listopada 2004, stwierdzenie 5.25.
Istnienie w randze osiem wynika z poprzedniego
twierdzenia po zmianie znaku formy.
Tutaj otrzymaliśmy jej wersję ujemnie określoną.
Parzystość oznacza, że wszystkie samoprzecięcia klas
całkowitych są parzyste, natomiast unimodularność
wiąże formę z homologiczną trywialnością brzegu.
Nie jest jednak prawdą, że każda całkowita sfera
homologiczna jest dyfeomorficzna z $S^3$.
\end{uwaga}

\begin{stwierdzenie}[Brzeg $E_8$ i sfera homologiczna Poincarégo]
\label{st:tg-E8-poincare}
Brzeg ujemnie określonego plumbingowania $P_{E_8}$
jest, z dokładnością do konwencji orientacyjnej,
\emph{sferą homologiczną Poincarégo}
$\Sigma(2,3,5)$. Jej grupa fundamentalna jest
nieprzemienna i nietrywialna; w szczególności
$\partial P_{E_8}$ nie jest homeomorficzny z $S^3$.
\end{stwierdzenie}

\begin{proof}[Identyfikacja przez rozwarstwienie Seiferta]
Graf $E_8$ jest gwiazdą z centralną wagą $-2$ i trzema
ramionami złożonymi odpowiednio z $1$, $2$ i $4$
wierzchołków o wadze $-2$. Ułamki łańcuchowe tych ramion
wynoszą $-2/1$, $-3/2$ i $-5/4$, co sprawdza się
indukcyjnie ze wzoru $[-2,\ldots,-2]^-=-(m+1)/m$.
Twierdzenie \ref{thm:tg-star-seifert},
przyjmowane tu wyprzedzająco z następnego rozdziału,
identyfikuje brzeg z rozmaitością o znormalizowanych
danych Seiferta
\[
M(0;-2;(2,1),(3,2),(5,4)).
\]
Identyfikację geometryczną plumbingu ze sferą
Brieskorna $\Sigma(2,3,5)$ przyjmujemy jako
odrębny wynik: R.~C. Kirby, M.~G. Scharlemann,
\emph{Eight faces of the Poincaré homology 3-sphere},
w: Geometric Topology, Academic Press, 1979,
opisy 1, 3, 4 i 6. W szczególności nie jest to
wniosek z samej zgodności homologii ani
z wartości wyznacznika.

Dla niebanalności grupy fundamentalnej nie trzeba
obliczać rzędu całej grupy. Stosujemy twierdzenie
\ref{thm:tg-seifert-presentation} z następnego rozdziału:
podstawa ma rodzaj zero, $b=-2$, a trzy pary
$(\alpha_i,\beta_i)$ to $(2,1),(3,2),(5,4)$.
Po oznaczeniu generatorów wyjątkowych włókien
przez $x_1,x_2,x_3$ uzyskujemy prezentację
\[
\pi_1(\partial P_{E_8})\cong
\big\langle x_1,x_2,x_3,h\mid
\begin{array}{l}
[h,x_i]=1\ (1\le i\le3),\\
x_1^2h=x_2^3h^2=x_3^5h^4=1,\\
x_1x_2x_3=h^{-2}
\end{array}
\big\rangle.
\]
Element $h$ reprezentuje regularne włókno.
Ponieważ jest centralny, $\langle h\rangle$ jest
podgrupą normalną. Dodanie relacji $h=1$ opisuje
iloraz przez tę podgrupę: po usunięciu generatora
$h$ z prezentacji pozostają dokładnie poniższe
relacje. Otrzymujemy grupę trójkątną
\[
\Delta(2,3,5)=\langle x_1,x_2,x_3\mid
x_1^2=x_2^3=x_3^5=x_1x_2x_3=1\rangle.
\]
Do wykazania nieprzemienności wystarczy
jawna surjekcja na $A_5$, bez używania pełnej
klasyfikacji grup trójkątnych.
Przy konwencji składania permutacji od prawej
niech $s=(12)(34)$ i $t=(235)$.
Wówczas $s$ ma rząd dwa, $t$ ma rząd trzy,
a $st=(1\,2\,4\,3\,5)$ ma rząd pięć.
Podgrupa $\langle s,t\rangle$ ma więc rząd
podzielny przez $2,3,5$ i dzielący $60$.
Gdyby miała rząd $30$, byłaby podgrupą
indeksu dwa w $A_5$. Taka podgrupa nie istnieje:
homomorfizm ilorazowy do $\mathbb Z/2$
musiałby zerować wszystkie 3-cykle, a te
generują $A_5$. Istotnie, każda parzysta
permutacja jest iloczynem par transpozycji;
para mająca wspólny element jest 3-cyklem
lub identycznością, a para rozłączna jest
iloczynem dwóch 3-cykli, np.
$(ab)(cd)=(acb)(acd)$.
Zatem $\langle s,t\rangle=A_5$.
Przypisanie $x_1\mapsto s$, $x_2\mapsto t$,
$x_3\mapsto(st)^{-1}$ spełnia wszystkie
relacje grupy trójkątnej i daje surjekcję.
W konsekwencji istnieje surjekcja
$\pi_1(\partial P_{E_8})\twoheadrightarrow A_5$.
Grupa $A_5$ jest nieprzemienna (na przykład
$st\ne ts$), więc surjekcja wyklucza przemienność
grupy fundamentalnej brzegu, a tym bardziej jej
trywialność. Z drugiej strony, z twierdzenia
Hurewicza w stopniu pierwszym abelianizacja
tej grupy jest izomorficzna z $H_1(\partial P_{E_8};\mathbb Z)$,
które na mocy twierdzenia \ref{thm:tg-E8-unimodular}
jest zerowe. Otrzymujemy przykład przestrzeni będącej
sferą homologiczną, ale nie homotopijną.
\end{proof}

\begin{uwaga}[Pełna grupa fundamentalna]
W istocie $\pi_1(\Sigma(2,3,5))$ jest binarną grupą
dwudziestościenną rzędu $120$, czyli pewnym centralnym
rozszerzeniem $A_5$ przez $\mathbb Z/2\mathbb Z$.
Wynika to z opisu $\Sigma(2,3,5)\cong S^3/I^*$
w przytoczonej pracy Kirby'ego i Scharlemanna
(opis 4): grupa $I^*$ działa przez lewe
mnożenie jednostkowych kwaternionów na $S^3$,
swobodnie, a $S^3$ jest jednospójna.
Jest ona przeciwobrazem grupy obrotów
dwudziestościanu w dwukrotnym nakryciu
$SU(2)\to SO(3)$, stąd jej rząd wynosi $120$.
Pełne obliczenie tej grupy nie jest potrzebne do
wykazania, że brzeg jest różny od $S^3$.
\end{uwaga}

\section{Rachunek Kirby'ego, przesuwanie uchwytów i wydmuchiwanie}
\label{sec:tg-plumbing-kirby}

\begin{stwierdzenie}[Przesunięcie uchwytu a kongruencja macierzy]
Niech $W$ będzie czterorozmaitością zbudowaną z $B^4$
przez dołączenie $r$ uchwytów indeksu dwa, o macierzy
przecięcia $Q$ w bazie klas powierzchni domkniętych.
Przesunięcie $i$-tego uchwytu po $j$-tym odpowiada
zastąpieniu $e_i$ przez $e_i\pm e_j$. Macierz formy
w nowej bazie ma postać $P^TQP$, gdzie
$P\in GL_r(\mathbb Z)$ realizuje taką zmianę bazy.
W szczególności wyznacznik oraz sygnatura formy
pozostają niezmienione.
\end{stwierdzenie}
\begin{proof}
Forma przecięcia jest formą dwuliniową na
$H_2(W;\mathbb Z)$. Jeśli kolumny macierzy $P$
reprezentują nowe wektory bazowe w starej bazie,
to ich wzajemne iloczyny skalarne tworzą macierz
$P^TQP$. Geometryczne działanie przesunięcia na klasach
homologii zostało uzasadnione w stwierdzeniu
\ref{st:kt-handle-slide}; dla znaku
$\varepsilon=\pm1$ mamy $P=I+\varepsilon E_{ji}$,
gdzie $i\ne j$ i $E_{ji}$ jest macierzą
elementarną z jedynką w pozycji $(j,i)$.
Macierz $P$ ma wyznacznik $\pm1$, a zatem
$\det(P^TQP)=(\det P)^2\det Q=\det Q$.
Z prawa bezwładności Sylvestera wynika niezmienniczość
liczby dodatnich i ujemnych wartości własnych.
\end{proof}

\begin{przyklad}[Jawny rachunek przesunięcia]
Niech $Q=\begin{psmallmatrix}a&\ell\\\ell&b\end{psmallmatrix}$.
Po zastąpieniu pierwszego wektora $e_1$ przez
$e_1+e_2$ otrzymujemy
\[
Q'=\begin{pmatrix}a+2\ell+b&\ell+b\\\ell+b&b\end{pmatrix},
\qquad \det Q'=ab-\ell^2=\det Q.
\]
Rachunek pokazuje dokładnie, skąd bierze się wyraz
$2\ell$ w nowym ramowaniu: pochodzi z obu iloczynów
mieszanych formy symetrycznej.
\end{przyklad}

\begin{stwierdzenie}[Algebraiczne usunięcie liścia o wadze $-1$]
Niech $\Gamma$ będzie ważonym drzewem i niech $v$
będzie liściem o wadze $-1$, połączonym tylko z
wierzchołkiem $u$ o wadze $b$. Wtedy macierz
$Q_\Gamma$ jest kongruentna nad $\mathbb Z$
z sumą ortogonalną $[-1]\oplus Q_{\Gamma'}$, gdzie
$\Gamma'$ otrzymujemy przez usunięcie liścia $v$
i zwiększenie wagi $u$ z $b$ do $b+1$.
W szczególności $\det Q_\Gamma=-\det Q_{\Gamma'}$.
\end{stwierdzenie}
\begin{proof}
Niech $e_v,e_u$ oznaczają odpowiednie wektory bazowe.
Dane przecięcia to $e_v^2=-1$, $e_v\cdot e_u=1$
i $e_v\cdot e_w=0$ dla $w\notin\{u,v\}$.
Wprowadźmy nową bazę: zostawiamy $e_v$, a $e_u$
zastępujemy przez $f_u=e_u+e_v$. Wtedy
$e_v\cdot f_u=1-1=0$ oraz
$f_u^2=b+2-1=b+1$. Dla pozostałych wierzchołków
$w$ mamy $f_u\cdot e_w=e_u\cdot e_w$, ponieważ
$e_v\cdot e_w=0$. Otrzymaliśmy więc dokładnie
wymaganą sumę ortogonalną. Zmiana bazy jest
unimodularna, ponieważ jej odwrotność uzyskujemy
przez $e_u=f_u-e_v$ przy zachowaniu pozostałych
wektorów. Równość wyznaczników wynika ze wzoru
na wyznacznik sumy blokowej. Ponadto
$\operatorname{coker}Q_\Gamma\cong
\operatorname{coker}Q_{\Gamma'}$,
ponieważ blok $[-1]$ ma trywialny kokernel.
W przypadku nieosobliwym forma splecenia jest
zachowana przez izomorfizm indukowany
kongruencją, jak w stwierdzeniu
\ref{st:kt-kirby-matrix}; blok $[-1]$ wnosi
jedynie wartość całkowitą, zerową modulo $\mathbb Z$.
\end{proof}

\begin{uwaga}[Geometria operacji blow-down]
Jeżeli klasę $e_v$ reprezentuje gładko zanurzona sfera
o samoprzecięciu $-1$, to jej otoczenie tubularne jest
ujemną wiązką dyskową $E(-1)$, której brzeg jest
$S^3$ (przykład rozkładu uchwytowego
$\overline{\mathbb{CP}}^{\,2}$ w rozdziale
\ref{ch:kt-cztery}).
Wybieramy sferę w obrębie wnętrza rozmaitości.
Można zastąpić jej otoczenie przez kulę
$B^4$, sklejając z zachowaniem orientacji
rozmaitości; nazywa się to \emph{wydmuchiwaniem w dół}
(\emph{blow-down}). Dla liścia plumbingowego operacja
prowadzi do grafu $\Gamma'$ ze zmienioną wagą:
w standardowym diagramie liść jest niezwęźlonym
meridianem składowej $u$ z ramowaniem $-1$.
Geometryczny ruch blow-down usuwa tę składową
i wprowadza pełny skręt pojedynczej nici;
po izotopii pozostaje ta sama krzywa
przyczepienia ze zwiększonym o jeden ramowaniem.
Jest to lokalny ruch rachunku Kirby'ego
omówiony w rozdziale \ref{ch:kt-cztery},
a jego wpływ na współczynnik chirurgii
zgadza się z regułą slam-dunk ze
stwierdzenia \ref{st:tg-plumbing-chain-surgery}.
Przeciwnie, \emph{wydmuchanie} (\emph{blow-up}) dodaje
lokalną sferę wyjątkową o samoprzecięciu $-1$.
Nie jest to dowolne skreślanie wiersza macierzy:
konieczna jest geometryczna obecność odpowiedniej sfery.
\end{uwaga}

\begin{center}
\begin{tikzpicture}[scale=.95,>=Stealth]
  \node[tgRedVertex,minimum size=11mm] (a) at (-2.5,0) {$-1$};
  \node[tgBlueVertex,minimum size=11mm] (b) at (-.8,0) {$b$};
  \draw[tgAccent] (a)--(b);
  \draw[tgTeal,-{Stealth[length=3.3mm]},line width=1.6pt] (.25,0)--(1.40,0);
  \node[tgTealVertex,minimum size=11mm] (c) at (2.65,0) {$b+1$};
  \node[text=tgNavy,font=\footnotesize] at (.05,-1.00)
    {Algebraiczne usunięcie liścia $-1$; w geometrii -- blow-down.};
\end{tikzpicture}
\end{center}

\section{Plumbing \texorpdfstring{$E_8$}{E8}, twierdzenia Freedmana i Donaldsona}
\label{sec:tg-plumbing-freedman}

\begin{twierdzenie}[Donaldson; diagonalizacja form określonych]
\label{thm:tg-donaldson-definite}
Niech $X$ będzie gładką, zamkniętą, spójną,
zorientowaną czterorozmaitością.
Jeśli forma przecięcia $Q_X$ na grupie
$H_2(X;\mathbb Z)/\operatorname{Tor}H_2(X;\mathbb Z)$
jest dodatnio lub ujemnie określona,
to jest izometryczna nad $\mathbb Z$ odpowiednio
z formą diagonalną
$\langle1\rangle^{\oplus b_2(X)}$ albo
$\langle-1\rangle^{\oplus b_2(X)}$.
\end{twierdzenie}

Wynik Simona Donaldsona jest głębokim twierdzeniem
teorii połączeń antysamodualnych. Przyjmujemy go bez dowodu.
Pierwotna wersja dla rozmaitości jednospójnych:
S.~K. Donaldson, \emph{An application of gauge theory
to four-dimensional topology}, J. Differential Geom.
18 (1983), 279--315, twierdzenie A.
Podana wyżej ogólna wersja:
S.~Friedl, M.~Nagel, P.~Orson, M.~Powell,
\emph{The foundations of four-manifold theory
in the topological category}, NYJM Monographs
6 (2025), twierdzenie 8.11 i uwaga 8.12.
Przypadek ujemnie określony wynika z dodatniego
po odwróceniu orientacji, które zmienia znak
formy przecięcia.

W szczególności żadna gładka zamknięta
zorientowana czterorozmaitość nie ma formy
przecięcia $Q_{E_8}$ ani $-Q_{E_8}$.
Parzystość jest niezmiennikiem całkowitej
izometrii, a formy te są parzyste, podczas gdy
$\pm I_8$ ma wektor bazowy o nieparzystym
samoprzecięciu $\pm1$.

\begin{twierdzenie}[Przeszkoda dla gładkich wypełnień homologicznych]
\label{thm:tg-E8-no-smooth-fill}
Niech $Y=\partial P_{E_8}$ będzie zorientowanym brzegiem
ujemnego plumbingu $E_8$. Żadna z orientacji $Y$
nie jest brzegiem zwartej, gładkiej, orientowanej
czterorozmaitości $C$ o homologii całkowitej kuli,
tj. $H_0(C;\mathbb Z)\cong\mathbb Z$ i
$H_i(C;\mathbb Z)=0$ dla $i>0$.
W szczególności $Y$ nie jest brzegiem gładkiej
rozmaitości kontraktywnej.
\end{twierdzenie}
\begin{proof}
Przypuśćmy najpierw, że istnieje taka $C$ z
dyfeomorfizmem zachowującym orientację
$\partial C\cong -Y$. Sklejamy ją z $P_{E_8}$
przez odpowiadający mu dyfeomorfizm brzegów
odwracający orientację:
$X=P_{E_8}\cup_Y C$.
Gładkie kołnierze brzegów pozwalają wyposażyć
otoczenie miejsca sklejenia w strukturę
gładkiego produktu $Y\times(-\varepsilon,\varepsilon)$.
Otrzymujemy gładką, zamkniętą, spójną,
orientowaną czterorozmaitość.

Pogrubiamy obie części w tym kołnierzu,
uzyskując otwarte pokrycie, którego części
retraktują się deformacyjnie na $P_{E_8}$ i $C$,
a przecięcie na $Y$.
Ciąg Mayera--Vietorisa zawiera fragment
\[
H_2(Y)\longrightarrow
H_2(P_{E_8})\oplus H_2(C)
\longrightarrow H_2(X)\longrightarrow H_1(Y),
\]
gdzie wszystkie grupy mają współczynniki
w $\mathbb Z$. Skrajne grupy są zerowe,
bo $Y$ jest całkowitą sferą homologiczną.
Ponieważ także $H_2(C)=0$, włożenie
$P_{E_8}\hookrightarrow X$ indukuje izomorfizm
drugich grup homologii.
Nie potrzebujemy założenia $\pi_1(C)=1$.
Gdy $C$ jest kontraktywna, van Kampen daje
dodatkowo $\pi_1(X)=1$, ponieważ obie
części są jednospójne, a przecięcie spójne.
Ponadto przecięcia klas reprezentowanych przez sfery
we wnętrzu $P_{E_8}$ nie ulegają zmianie przy sklejeniu.
Stąd forma przecięcia zamkniętej $X$ jest dokładnie
$Q_{E_8}$: parzysta, unimodularna i ujemnie określona.
Przeczy to przytoczonemu twierdzeniu Donaldsona,
gdyż $Q_{E_8}$ nie jest diagonalizowalna nad $\mathbb Z$.
Dowiedliśmy nieistnienia gładkiej całkowitej
kuli homologicznej o brzegu $-Y$.
Gdyby natomiast istniało takie wypełnienie
brzegu $Y$, to po odwróceniu orientacji
czterorozmaitości otrzymalibyśmy gładką
całkowitą kulę homologiczną o brzegu $-Y$,
co również prowadziłoby do sprzeczności.
Zatem przeszkoda dotyczy obu orientacji
sfery homologicznej Poincarégo.
\end{proof}

\begin{twierdzenie}[Freedman; topologiczne wypełnianie sfer homologicznych]
\label{thm:tg-freedman-homology-fill}
Każda zamknięta, zorientowana całkowita sfera homologiczna
trójwymiarowa $Y$ jest brzegiem zwartej,
kontraktywnej \emph{topologicznej} czterorozmaitości $C$.
Orientację $C$ wybieramy tak, by identyfikacja
$\partial C\cong Y$ zachowywała orientację.
Nie zakłada się, że $C$ ma strukturę gładką.
\end{twierdzenie}

Przyjmujemy to twierdzenie z literatury:
M.~H. Freedman, F.~Quinn,
\emph{Topology of 4-Manifolds}, Princeton, 1990,
wniosek 9.3C, s.~146. Jego dowód
korzysta z czterowymiarowej topologicznej teorii chirurgii
i twierdzenia o zanurzaniu dysków. Właśnie różnica między
słowami \emph{topologiczna} i \emph{gładka} jest tu istotą
zagadnienia, a nie jedynie techniczną uwagą.

\begin{wniosek}[Zamknięta topologiczna rozmaitość $E_8$]
\label{cor:tg-closed-E8}
Istnieje zamknięta, spójna, zorientowana, po prostu
spójna topologiczna czterorozmaitość $X_{E_8}$,
której całkowita forma przecięcia jest równa
ujemnej formie $E_8$. Nie dopuszcza ona struktury gładkiej.
\end{wniosek}
\begin{proof}
Niech $Y=\partial P_{E_8}$. Z udowodnionego rachunku
homologii wynika, że $Y$ jest całkowitą sferą homologiczną.
Z twierdzenia \ref{thm:tg-freedman-homology-fill}
istnieje kontraktywna topologiczna czterorozmaitość
$C$ o brzegu $-Y$. Sklejamy
$X_{E_8}=P_{E_8}\cup_Y C$ przez homeomorfizm
brzegów odwracający orientację.
Twierdzenie o kołnierzu dla rozmaitości
topologicznych zapewnia lokalny model produktu
także w kategorii topologicznej. Otrzymana
przestrzeń jest zatem zamkniętą spójną
orientowaną czterorozmaitością topologiczną.

W otwartym pokryciu po pogrubieniu kołnierzy
obie części są jednospójne, a ich przecięcie
spójne. Z twierdzenia van Kampena wynika
$\pi_1(X_{E_8})=1$. Rachunek Mayera--Vietorisa
z dowodu twierdzenia \ref{thm:tg-E8-no-smooth-fill}
daje izomorfizm $H_2(P_{E_8};\mathbb Z)\to
H_2(X_{E_8};\mathbb Z)$ indukowany włożeniem.
Sekcje zerowe leżą w gładkim wnętrzu
$P_{E_8}$ i ich lokalne liczby przecięć
pozostają niezmienione.
Równoważnie można obliczyć formę przecięcia
przez dualność Poincarégo i iloczyn kohomologiczny:
naturalność oraz wycinanie identyfikują
przecięcia klas o nośnikach w tym wnętrzu
z ich przecięciami w całej rozmaitości.
Stąd forma $X_{E_8}$ jest równa $Q_{E_8}$.
Ponadto $H_1(X_{E_8})=0$ z jednospójności,
$H_3(X_{E_8})=0$ z dualności i uniwersalnych
współczynników, a $H_0$ i $H_4$ są izomorficzne
z $\mathbb Z$. Zatem $b_2=8$ i $\chi(X_{E_8})=10$.
Gdyby $X_{E_8}$ dopuszczała strukturę gładką,
twierdzenie \ref{thm:tg-donaldson-definite}
wymuszałoby całkowitą diagonalizację jej formy
przecięcia, co jest niemożliwe ze względu
na parzystość $Q_{E_8}$.
\end{proof}

\begin{uwaga}[Niezależna przeszkoda Rochlina]
Twierdzenie Rochlina mówi, że sygnatura
zamkniętej gładkiej orientowanej rozmaitości
spin wymiaru cztery jest podzielna przez $16$.
W tym przykładzie gdyby istniała struktura gładka
na $X_{E_8}$, byłaby ona spin: $H_1(X_{E_8})=0$
i $H_2(X_{E_8})$ jest wolna. Uniwersalne
twierdzenie o współczynnikach identyfikuje
kohomologię stopnia dwa nad $\mathbb Z$
i nad $\mathbb F_2$ odpowiednio z
$\operatorname{Hom}(H_2(X_{E_8}),\mathbb Z)$
i $\operatorname{Hom}(H_2(X_{E_8}),\mathbb F_2)$,
gdyż składniki $\operatorname{Ext}$ z $H_1$
znikają. Przez naturalność redukcja
współczynników odpowiada redukcji wartości
tych homomorfizmów. Wolność $H_2$ zapewnia
jej surjektywność. Każda klasa modulo dwa
jest więc redukcją klasy całkowitej. Parzystość formy oznacza
zanikanie kwadratów tych klas modulo dwa.
Wzór Wu, przy orientowalności, daje
$\langle w_2\smile a,[X_{E_8}]_2\rangle=0$
dla każdego $a\in H^2(X_{E_8};\mathbb F_2)$.
Niezdegenerowanie dualności Poincarégo nad
$\mathbb F_2$ implikuje $w_2=0$, czyli
istnienie struktury spin.
Ponieważ $\operatorname{sign}(Q_{E_8})=-8$
nie jest podzielna przez $16$, otrzymujemy
drugą sprzeczność.
Przyjmujemy twierdzenie Rochlina i wzór Wu
z literatury: Friedl--Nagel--Orson--Powell,
cytowana monografia, twierdzenie 8.8
i uwaga 8.9. Nie należy w ogólności
wnioskować o istnieniu struktury spin
z samej parzystości formy na części wolnej
homologii bez kontroli torsji.
\end{uwaga}

\begin{uwaga}[Nie należy utożsamiać trzech rozmaitości]
Występują tu trzy różne obiekty:
\begin{enumerate}[label=(\roman*)]
\item $P_{E_8}$ jest \emph{gładką zwartą} czterorozmaitością
z niepustym brzegiem $\Sigma(2,3,5)$;
\item $C$ jest \emph{topologiczną kontraktywną}
czterorozmaitością z przeciwnie zorientowanym brzegiem;
\item $X_{E_8}=P_{E_8}\cup C$ jest \emph{zamkniętą
czterorozmaitością topologiczną}, która nie może być wygładzona.
\end{enumerate}
Sama obecność formy $E_8$ na rozmaitości \emph{z brzegiem}
nie przeczy twierdzeniu Donaldsona. Sprzeczność powstałaby
po gładkim domknięciu przez całkowitą kulę homologiczną,
ponieważ zachowałoby ono formę przecięcia $Q_{E_8}$.
Nie oznacza to zakazu dowolnego gładkiego domknięcia.
Na przykład podwojenie
$P_{E_8}\cup_Y(-P_{E_8})$ jest gładką zamkniętą
jednospójną czterorozmaitością: stosujemy kołnierze
i twierdzenie van Kampena jak powyżej.
Z ciągu Mayera--Vietorisa oraz
$H_1(Y)=H_2(Y)=0$ wynika, że włożenia obu części
dają izomorfizm drugiej grupy homologii z sumą
prostą ich drugich grup homologii.
Klasy z różnych części mają rozłączne reprezentanty
we wnętrzach, więc forma podwojenia to
$Q_{E_8}\oplus(-Q_{E_8})$.
Jest ona nieokreślona, a jej sygnatura wynosi zero;
twierdzenie Donaldsona o formach określonych
nie daje tu sprzeczności.
\end{uwaga}

\section{Odczytywanie topologii z grafu: zestawienie i dalsze kierunki}
\label{sec:tg-plumbing-conclusions}

\begin{przyklad}[Trzy grafy i trzy różne informacje]
\begin{enumerate}[label=(\roman*)]
\item Łańcuch $A_3$ z wagami $-2$ ma macierz
$Q=\begin{psmallmatrix}-2&1&0\\1&-2&1\\0&1&-2\end{psmallmatrix}$,
wyznacznik $-4$ i brzeg o grupie $H_1\cong\mathbb Z/4\mathbb Z$.
\item Gwiazda $D_4$ ma wyznacznik $4$, ale
$H_1\cong\mathbb Z/2\oplus\mathbb Z/2$.
\item Graf $E_8$ ma wyznacznik $1$, więc jego brzeg jest
sferą homologiczną, lecz nie jest sferą $S^3$.
\end{enumerate}
Wyznacznik dostarcza pierwszej informacji o brzegu,
postać Smitha rozróżnia możliwe struktury grupy,
a grupa fundamentalna i geometryczne dane sklejania
potrafią rozróżnić także rozmaitości o identycznej homologii.
\end{przyklad}

\begin{uwaga}[Zakres stosowalności metody]
Wielkości $\det Q$, $\operatorname{sign}Q$ i
$\operatorname{coker}Q$ są dostępne przez elementarne rachunki
na macierzach. Nie wyznaczają jednak samodzielnie
typu homeomorfizmu ani dyfeomorfizmu brzegu.
Różne ważone grafy mogą opisywać tę samą rozmaitość,
a ich porównanie wymaga operacji Kirby'ego i odpowiednich
ruchów na grafach. Z kolei przy konstrukcji z pętlami,
powierzchniami bazowymi o dodatnim rodzaju czy grafami
o znakowanych krawędziach trzeba zachować dodatkowe
dane geometryczne: sam wektor wag nie wystarcza.
\end{uwaga}

\begin{uwaga}[Dalsze zastosowania]
Grafy plumbingowe pojawiają się w rozdzielczościach
osobliwości powierzchni zespolonych, teorii splotów
algebraicznych, rozmaitościach Seiferta, rachunku Kirby'ego,
a także w badaniach niezmienników Floera i nierówności
dla gładkich wypełnień. Bardziej zaawansowana konstrukcja
\emph{rational blow-down} zastępuje odpowiednie ujemnie
określone plumbingowanie przez racjonalną kulę homologiczną;
wymaga ona dodatkowej wiedzy o gładkich czterorozmaitościach.
Tutaj ograniczamy się do fundamentów potrzebnych do
rozumienia form przecięcia i przypadku $E_8$.
\end{uwaga}

\medskip
\noindent\textbf{Literatura do rozdziału.}
R.~E. Gompf, A.~I. Stipsicz,
\emph{4-Manifolds and Kirby Calculus}, AMS, 1999;
M.~H. Freedman, F.~Quinn,
\emph{Topology of 4-Manifolds}, Princeton, 1990;
S.~K. Donaldson, \emph{An application of gauge theory
to four-dimensional topology},
J. Differential Geom. 18 (1983), 279--315;
S.~Friedl, M.~Nagel, P.~Orson, M.~Powell,
\emph{The foundations of four-manifold theory
in the topological category}, NYJM Monographs 6 (2025);
W.~D. Neumann, \emph{A calculus for plumbing applied to the
topology of complex surface singularities and degenerating complex
curves}, Trans. Amer. Math. Soc. 268 (1981), 299--344;
R.~C. Kirby, \emph{The Topology of 4-Manifolds},
Lecture Notes in Math. 1374, Springer, 1989.


% Iteracja 8: uzupelnienie teorii Seiferta i rachunku Kirby'ego.
